\section{A critical primorial clock}

Let
\[
 P_L=\prod_{\substack{\ell\le L\\ \ell\ \mathrm{prime}}}\ell,
 \qquad
 j_L=\left\lfloor\frac{L\log L-\log P_L}{\log2}\right\rfloor,
 \qquad Q_L=2^{j_L}P_L . \tag{1}
\]
The prime number theorem in the form $\log P_L\sim L$ makes $j_L$ nonnegative for all
large $L$.  The floor in (1) gives
\[
 \log Q_L=L\log L+O(1),
 \qquad
 L\sim\frac{\log Q_L}{\log\log Q_L}. \tag{2}
\]
Every prime divisor of $Q_L$ is at most $L$.

\begin{lemma}[Endpoint shell rows]\label{lem:endpoints}
For all sufficiently large $L$, there are primes $p_L,r_L$ satisfying
\[
 Q_L<p_L<Q_L+\frac{Q_L}{2L},
 \qquad
 2Q_L-\frac{Q_L}{2L}<r_L<2Q_L. \tag{3}
\]
Consequently $R_{p_L}=L$ and $R_{r_L}=2L-1$.
\end{lemma}

\begin{proof}
The classical zero-free-region remainder gives, for some $c>0$,
\[
 \pi(x)=\operatorname{Li}(x)+O\!\left(xe^{-c\sqrt{\log x}}\right);
\]
see, for example, \citet{davenport2000,montgomeryvaughan2007} and
\citet{iwanieckowalski2004}.  Set $x=Q_L$ and $y=Q_L/(2L)$.  The main term in
$\pi(x+y)-\pi(x)$ is asymptotic to $y/\log x$, whereas the error-to-main ratio is
\[
 O\!\left(L\log Q_L\,e^{-c\sqrt{\log Q_L}}\right)=o(1)
\]
by (2).  This gives the first prime in (3); the interval immediately below $2Q_L$ is
identical.  Multiplying (3) by $L/Q_L$ yields values strictly between $L$ and
$L+1/2$, and between $2L-1/2$ and $2L$, respectively.  Taking floors proves the
cutoff identities.
\end{proof}

\begin{theorem}[Critical raw-mass separation]\label{thm:main}
Along the clocks (1),
\[
 N_{p_L}=2,
 \qquad
 N_{r_L}=2\{1+\pi(2L-1)-\pi(L)\}, \tag{4}
\]
and therefore
\[
 \kappa_{\raw}(Q_L,L)
 \ge1+\pi(2L-1)-\pi(L)
 \sim\frac{L}{\log L}\longrightarrow\infty. \tag{5}
\]
\end{theorem}

\begin{proof}
Put $h_L=4LQ_L$.  Every $2\le m\le L$ has a prime divisor at most $L$, hence a
divisor of $Q_L$.  Thus $A_{h_L}(L)=1$.

If $L<m<2L$ is composite, it has a prime divisor at most
$\sqrt m<\sqrt{2L}<L$ for $L\ge3$, so it is not primitive modulo $h_L$.  If $m$ is
prime, then $m>L$ and it divides neither $Q_L$ nor $4L$, so it is primitive.  Hence
the positive primitive multipliers up to $2L-1$ are exactly $1$ and the primes in
$(L,2L)$.  Lemma~\ref{lem:endpoints} now gives (4).  The PNT on $(L,2L)$ gives (5).
\end{proof}

Theorem~\ref{thm:main} refutes a scale-independent raw comparability constant as a
theorem of the modeled support.  It does not refute comparability for a particular
source after nonuniform weighting or normalization.
